\subsection{Divisorial components of threshold subsolutions}

The linear Gauduchon threshold of \cref{cor:gauduchon-threshold} controls how much a divisorial component lowers the coefficient in a Monge--Amp\`ere inequality. For this subsection, fix a Hermitian form $\omega$ on a connected compact complex manifold of dimension $n\geq2$, a smooth probability volume $\mathrm{d}V$, and a nonzero $0\leq f\in L^1(\mathrm{d}V)$. Write $p=f^{1/n}$ and $b=b_\omega(f)$.

\begin{proposition}\label{prop:threshold-divisorial-constraint}
Suppose $u\in\PSH(X,\omega)$ satisfies
\[
 \mathrm{MA}^{\mathrm{ext}}_\omega(u)\geq c f\,\mathrm{d}V,
 \qquad \omega+\ddc u\geq[E],
\]
where $c>0$ and $E$ is an effective real divisor. Let $\beta_E$ be a smooth closed representative of its Bott--Chern class, and put $\kappa=\kappa_\omega(\beta_E)$. Then $b>0$, $0\leq\kappa<1$ and
\begin{equation}\label{eq:divisorial-coefficient-loss}
 c\leq b(1-\kappa)^n.
\end{equation}
If $c=b$, every smooth maximizing sequence $(T_j)$ for $D_p$ satisfies
\[
 \int_X\beta_E\wedge T_j=\int_X[E]\wedge T_j\longrightarrow0.
\]
Every weak cluster current $T$ therefore satisfies $\int_X\beta_E\wedge T=0$, and the class of $E$ is not big. If also $p\in\mathcal D_\omega^0$, each such cluster current maximizes $D_p$.
\end{proposition}
\begin{proof}
The subsolution characterization gives $b\geq c>0$. For a smooth strictly positive normalized Gauduchon form $T$, \cref{lem:unbounded-gauduchon-trace-density} shows that
\[
 R=(\omega+\ddc u)\wedge T-c^{1/n}p k_T\,\mathrm{d}V
\]
is positive, with mass $1-c^{1/n}D_p(T)$. The density term gives no mass to the support of $E$, whereas $[E]\wedge T$ is supported there. Restricting the positive measure $R$ to that support and using $\omega+\ddc u\geq[E]$ gives
\[
 \int_X[E]\wedge T\leq R(X)=1-c^{1/n}D_p(T).
\]
Write $[E]=\beta_E+\ddc\ell_E$ by the Poincar\'e--Lelong formula. Since $\ell_E$ is locally integrable and $T$ is smooth and $\ddc$-closed, distributional Stokes identifies the left side with $\int_X\beta_E\wedge T$. It is nonnegative for every such test, so $\kappa\geq0$. Also $D_p(T)>0$, which gives $\kappa<1$. Taking the supremum of $D_p(T)$ and using $\rho_\omega(p)=b^{-1/n}$ proves \eqref{eq:divisorial-coefficient-loss}.

If $c=b$ and $(T_j)$ is maximizing, the right side of the preceding measure bound tends to zero. Pairing with the fixed smooth form $\beta_E$ is continuous under current convergence, proving the cluster assertion. Equation~\eqref{eq:divisorial-coefficient-loss} forces $\kappa=0$, so \cref{cor:gauduchon-threshold} excludes bigness. Under the tail hypothesis, the near-maximizer statement in \cref{prop:integrable-threshold-adjoints} makes every cluster a maximizing current.
\end{proof}

At the threshold $c=b$, if $E\ne0$, the vanishing pairing defines a nonempty proper exposed face of the weak closure of smooth strictly positive normalized Gauduchon forms. Indeed, this closure is compact and convex; the pairing with $\beta_E$ is continuous and nonnegative there, and it is strictly positive on every smooth strictly positive member. The proposition places every cluster of a smooth maximizing sequence on this face. This assertion concerns that closure; it does not assert that every normalized Gauduchon current is a limit of smooth ones. Nor does the smooth pairing $\int_X\beta_E\wedge T$ define $[E]\wedge T$ for a singular $T$.

The hypotheses of the next two corollaries on $b$ and on the maximizing sequence hold whenever $p\in\mathcal D_\omega^0$, since this gives $b>0$ and uniform integrability over the whole of $\mathscr G_\omega$ (\cref{thm:integrable-current-attainment}).

\Needspace{12\baselineskip}
\begin{corollary}\label{cor:semipositive-divisor-exclusion}
Assume $b>0$ and that $D_p$ has a smooth maximizing sequence $(T_j)$ with uniformly integrable densities $p k_{T_j}$. Let $E$ be an effective real divisor whose class has a smooth semipositive representative $\beta_E$. If
\[
 \int_{\{\beta_E=0\}}f\,\mathrm{d}V=0,
\]
then no threshold subsolution $u$ for $f$ satisfies $\omega+\ddc u\geq[E]$.
\end{corollary}
\begin{proof}
Suppose such a subsolution exists. By \cref{prop:threshold-divisorial-constraint}, there is a maximizing current $T\in\mathscr G_\omega$ with $\int_X\beta_E\wedge T=0$. We use the matrix-density notation of \cref{sec:weak-gauduchon-duality}: $\sigma_T=\omega\wedge T$ is the trace measure of $T$, $P_T$ is its matrix density and $B$ is the coefficient matrix of $\beta_E$. Since both matrices are positive semidefinite, the vanishing pairing gives $\operatorname{tr}(BP_T)=0$ at $\sigma_T$-almost every point. Wherever $B\ne0$, this forces $P_T$ to be singular: a positive definite $P_T$ would have strictly positive trace against the nonzero positive semidefinite matrix $B$. The determinant measure $\Gamma_0(T)$ vanishes where $P_T$ is singular, so it is concentrated on $\{\beta_E=0\}$. By \cref{thm:gaud-determinant-subcritical}, $\Gamma_0(T)=k_T\,\mathrm{d}V$, and the hypothesis on $f$ gives $p k_T=0$ almost everywhere. This contradicts $D_p(T)=\rho_\omega(p)>0$.
\end{proof}


\begin{corollary}\label{cor:nef-divisor-threshold-volume}
Let $0\leq f\in L^1(\mathrm{d}V)$ be nonzero, with $f^{1/n}\in\mathcal D_\omega^0$. Suppose a subsolution at $b_\omega(f)$ satisfies $\omega+\ddc u\geq[E]$ for an effective real divisor $E$ whose Bott--Chern class is nef. Then
\[
 \lambda_\mu(\beta_E)=0
\]
for every smooth positive probability measure $\mu$, where $\beta_E$ is any smooth closed representative of that class.
\end{corollary}
\begin{proof}
By \cref{prop:threshold-divisorial-constraint}, there is a maximizing current $T\in\mathscr G_\omega$ with $\int_X\beta_E\wedge T=0$. If $\lambda_\mu(\beta_E)>0$, \cref{prop:nef-annihilator-determinant} gives $\Gamma_0(T)=0$. It follows that $D_{f^{1/n}}(T)=0$, contradicting its positive maximal value. Since nefness gives $\lambda_\mu(\beta_E)\geq0$, the conclusion follows for each $\mu$.
\end{proof}

This criterion uses nefness and the variational volume, without requiring a smooth semipositive representative. It complements \cref{cor:semipositive-divisor-exclusion}: the Hopf hyperplane class has zero variational volume, yet its nowhere vanishing semipositive representative excludes a component at the threshold. Neither argument identifies positive nef volume with bigness or classifies all possible divisorial components.


These are necessary conditions on divisorial singularities at the maximal global coefficient. They do not apply unchanged to a smaller coefficient or to a threshold defined within a prescribed singularity type. Determining which annihilating currents occur, and whether these conditions can be supplemented to characterize the possible divisorial components, remains a geometric realization problem. No identification with a null locus or a non-K\"ahler locus is asserted.
